\documentclass[aspectratio=169,11pt,professionalfonts]{beamer}
% Paper 3 presentation -- all 28 figures from presentation/build_figures.py.
% Compiles with placeholders if the figure PDFs have not yet been rendered.
\usepackage[T1]{fontenc}
\usepackage[utf8]{inputenc}
\usepackage[spanish,es-nodecimaldot]{babel}
\usepackage{lmodern}
\usepackage{amsmath,amssymb,booktabs,graphicx}
\usepackage{xcolor}
\usepackage{hyperref}

\definecolor{ink}{HTML}{17283D}
\definecolor{accentblue}{HTML}{0072B2}
\definecolor{accentgreen}{HTML}{009E73}
\definecolor{accentorange}{HTML}{D55E00}
\definecolor{slate}{HTML}{68798A}
\definecolor{pale}{HTML}{EFF5F9}
\definecolor{lightborder}{HTML}{DDE6EC}

\setbeamertemplate{navigation symbols}{}
\setbeamertemplate{itemize items}[circle]
\setbeamercolor{normal text}{fg=ink,bg=white}
\setbeamercolor{frametitle}{fg=ink,bg=white}
\setbeamerfont{frametitle}{size=\large,series=\bfseries}
\setbeamercolor{structure}{fg=accentblue}
\setbeamercolor{block title}{fg=white,bg=accentblue}
\setbeamercolor{block body}{fg=ink,bg=pale}
\setbeamercolor{keybox}{fg=ink,bg=pale}
\setbeamercolor{missingfigure}{fg=slate,bg=pale}
\setbeamercolor{footline}{fg=slate,bg=white}
\setbeamertemplate{headline}{}
\setbeamertemplate{footline}{%
  \leavevmode\hbox{%
  \begin{beamercolorbox}[wd=.76\paperwidth,ht=2.9ex,dp=1.25ex,leftskip=.45cm]{footline}%
  \tiny MÉTODOS NUMÉRICOS \quad\textcolor{accentblue}{\textbullet}\quad
  \insertsectionhead
  \end{beamercolorbox}%
  \begin{beamercolorbox}[wd=.24\paperwidth,ht=2.9ex,dp=1.25ex,rightskip=.45cm plus 1fil]{footline}%
  \hfill\scriptsize\insertframenumber\,/\,\inserttotalframenumber
  \end{beamercolorbox}}%
}
\setbeamertemplate{frametitle}{%
  \vskip.18cm\hskip.44cm\insertframetitle\par
  \vspace{.11cm}\hskip.44cm{\color{lightborder}\rule{\dimexpr\paperwidth-2.10cm\relax}{.65pt}}\vspace{-.14cm}%
}
\setlength{\parskip}{.4em}
\newcommand{\Sm}{S}
\newcommand{\lam}{\lambda}
\newcommand{\CVO}{F_{d,L,h}}
\newcommand{\SourceLine}[1]{\par\vspace{.05em}{\scriptsize\color{slate}#1\par}}
\newcommand{\FigureSlide}[4]{%
\begin{frame}[t]{#2}
  \vspace{-.14cm}
  \begin{center}
  \IfFileExists{figures/#1.pdf}{%
    \includegraphics[width=.965\linewidth,height=.515\textheight,keepaspectratio]{figures/#1.pdf}%
  }{%
    \fcolorbox{lightborder}{pale}{%
      \parbox[c][.51\textheight][c]{.89\linewidth}{%
        \centering{\large\bfseries Figura pendiente de generar}\par
        \vspace{.35em}{\scriptsize\texttt{\detokenize{#1.pdf}}}\par
        \vspace{.35em}{\scriptsize Genera los 28 PDF con el script de figuras.}%
      }%
    }%
  }
  \end{center}
  \vspace{-.16cm}
  \begin{beamercolorbox}[wd=.97\linewidth,sep=.32em,rounded=true]{keybox}
    \footnotesize\textbf{\textcolor{accentblue}{Interpretación.}} #3
  \end{beamercolorbox}
  \par\vspace{.13cm}{\tiny\color{slate}#4\par}
\end{frame}
}
\newcommand{\Act}[3]{%
\begin{frame}[plain]
  \vspace{1.18cm}
  {\color{accentblue}\Large\bfseries #1}\par
  \vspace{.27cm}
  {\color{ink}\Huge\bfseries #2}\par
  \vspace{.47cm}
  {\color{slate}\large #3}\par
  \vfill
  {\color{accentblue}\rule{3.0cm}{3pt}}
\end{frame}
}
\title{Recuperación numérica de mínimos multimodales\\de suavizamiento óptimo para pronóstico}
\subtitle{Búsqueda adaptativa, extremos exactos y aislamiento de Sturm}
\author{Heriberto Espino Montelongo}
\institute{Universidad de las Am\'ericas Puebla}
\date{Presentación de investigación \quad \textbar\quad octubre de 2026}
\begin{document}

\begin{frame}[plain]
  \vspace{.55cm}
  {\color{accentblue}\small PRESENTACIÓN DE INVESTIGACIÓN \quad | \quad OPTIMIZACIÓN NUMÉRICA}\par
  \vspace{.55cm}
  {\Huge\bfseries Recuperación numérica de mínimos multimodales\\de suavizamiento óptimo para pronóstico}\par
  \vspace{.36cm}
  {\Large\color{slate}Búsqueda adaptativa, extremos exactos y aislamiento de Sturm}\par
  \vfill
  {\large Heriberto Espino Montelongo}\par
  Universidad de las Am\'ericas Puebla\par
  {\color{slate}octubre de 2026 \quad\textbullet\quad Basada en el documento de trabajo 3}
  \vspace{.35cm}
\end{frame}

\section{La pregunta numérica}
\begin{frame}{¿Qué problema resuelve este trabajo?}
\begin{block}{El problema computacional}
Para un orden de diferencias $d$, ventana $L$ y horizonte $h$ fijos,
buscamos y comparamos los mínimos pertinentes de la pérdida de validación
de pronóstico $F_{d,L,h}(S)$ sobre $S\in[0,1]$.
\end{block}
\begin{columns}[T,totalwidth=\linewidth]
\begin{column}{.49\linewidth}
\textbf{Elementos ya conocidos}
\begin{itemize}
\item Mínimos cuadrados penalizados con diferencias finitas.
\item Pérdida cronológica de validación de pronósticos.
\item Optimización escalar y búsqueda de raíces.
\end{itemize}
\end{column}
\begin{column}{.49\linewidth}
\textbf{Nuestra contribución numérica}
\begin{itemize}
\item Geometría multimodal de la pérdida.
\item Tratamiento exacto de los extremos.
\item Detección adaptativa de puntos estacionarios.
\item Aislamiento de raíces para objetivos racionales.
\end{itemize}
\end{column}
\end{columns}
\SourceLine{No proponemos un nuevo suavizador PLS ni la primera CV predictiva; tampoco afirmamos superioridad predictiva universal.}
\end{frame}

\begin{frame}{Modelo y función objetivo: alcance de los resultados}
\[
\widehat{\tau}_{\lambda} =
\underset{\tau}{\arg\min}\left\{\lVert y-\tau\rVert_2^2+
\lambda\lVert D_d\tau\rVert_2^2\right\}
=H_\lambda y,\qquad
H_\lambda=(I+\lambda D_d^\top D_d)^{-1}.
\]
\[
F_{d,L,h}(S)=\frac{1}{|\mathcal O|}
\sum_{T\in\mathcal O}\frac{1}{h}
\left\|z_T-G_{d,h}H_{\lambda(S)}x_T\right\|_2^2.
\]
\begin{block}{Supuestos}
Ponderación identidad $V=I$, media de diferencias penalizadas $\mu=0$,
observaciones equidistantes, $(d,L,h)$ fijos y extrapolación natural de orden $d$.
\end{block}
\SourceLine{Cada origen utiliza solo los datos disponibles hasta ese instante; los bloques futuros realizados definen el error.}
\end{frame}

\section{La función objetivo}
\Act{PARTE I}{La función objetivo}{Suavizamiento, pronóstico y transformación del dominio.}
\FigureSlide{01_rolling_cv_definition}{Validación cruzada cronológica}{La pérdida promedia errores de pronóstico a $h$ pasos desde orígenes históricos ya completados. Cada origen utiliza únicamente observaciones anteriores.}{Esquema y curvas ilustrativos; no son resultados del experimento principal.}

\FigureSlide{02_smoothing_family}{Cómo cambia la tendencia con el suavizamiento}{Con $S$ pequeño se conservan más fluctuaciones; con $S$ grande se atenúan las diferencias penalizadas. Más suavizamiento no garantiza mejores pronósticos.}{Serie pedagógica con semilla fija; no es el experimento de Monte Carlo.}

\FigureSlide{03_lambda_to_s}{De una penalización ilimitada a un índice compacto}{La transformación creciente $S(\lambda)$ lleva $\lambda\in[0,\infty]$ a $S\in[0,1]$. Cerca del extremo superior se requiere especial resolución numérica.}{Espectro de diferencias finitas para $L=63$, $d=2$.}

\FigureSlide{04_effective_degrees_of_freedom}{Grados de libertad efectivos y suavizamiento}{Para el suavizador PLS, $\mathrm{EDF}=\mathrm{tr}(H_\lambda)=L-(L-d)S$. Las $d$ direcciones no penalizadas permanecen cuando $S\to1$.}{Identidad exacta con media de diferencias cero y $\lambda$ fijo.}

\FigureSlide{05_exact_endpoints}{Los extremos representan estimadores distintos}{En $S=0$ tenemos $H=I$; en $S=1$ obtenemos la proyección sobre $\ker D_d$. Este último caso no se sustituye por una $\lambda$ finita muy grande.}{Serie ilustrativa y extrapolación natural con $d=2$.}

\FigureSlide{06_loss_pooling}{Por qué combinar errores de distintos orígenes}{Cada origen histórico puede favorecer un grado de suavizamiento diferente. Optimizamos el promedio de pérdidas ya observadas.}{Tres ejemplos pequeños construidos con PLS; solo explicación.}

\section{Geometría y algoritmos}
\Act{PARTE II}{Geometría y algoritmos}{Derivadas, búsqueda adaptativa y aislamiento exacto en ejemplos pequeños.}
\FigureSlide{07_objective_derivative_curvature}{Pérdida, primera derivada y curvatura}{Las raíces de $F'(S)$ señalan puntos estacionarios. La curvatura clasifica extremos simples, pero los puntos planos requieren más cuidado.}{Ejemplo racional de seis observaciones previamente documentado.}

\FigureSlide{08_adaptive_search_steps}{Algoritmo adaptativo paso a paso}{Primero se usa una malla gruesa; después se refuerzan los extremos, se subdividen intervalos y Brent refina raíces detectadas. Los extremos exactos siguen siendo candidatos.}{Etapas construidas con las muestras finales; no es un registro cronológico del algoritmo.}

\FigureSlide{09_endpoint_refinement_ablation}{Qué falla sin refinar cerca de los extremos}{En 216 superficies, eliminar el refinamiento de frontera produjo cinco mínimos no detectados y tres casos con arrepentimiento positivo.}{Datos congelados de sensibilidad; la malla densa es una referencia aproximada.}

\FigureSlide{10_sturm_worked_example}{Conteo exacto de raíces en una instancia racional}{Para $L=6$, $d=2$, $h=2$, el numerador de la derivada tiene tres raíces positivas: mínimo, máximo y mínimo. Las posiciones dibujadas son aproximadas.}{Conteo exacto de Sturm $5-2=3$ para esta instancia pequeña.}

\FigureSlide{11_lambda_vs_s_landscape}{Un mismo objetivo, dos coordenadas}{La transformación monótona $S(\lambda)$ conserva los puntos estacionarios interiores y su clasificación, pero modifica la distancia numérica entre ellos.}{Mismo problema racional de seis observaciones en ambas coordenadas.}

\FigureSlide{12_adversarial_gallery}{Qué geometrías complican la búsqueda}{Los casos analíticos incluyen raíces cercanas, mínimos de frontera, mínimos planos e inflexiones estacionarias.}{Objetivos analíticos con soluciones conocidas; no necesariamente son pérdidas de PLS.}

\FigureSlide{13_adaptive_sampling_density}{Dónde concentra evaluaciones el algoritmo}{El método distribuye las evaluaciones de manera no uniforme, según la estructura de las derivadas. Esta visualización corresponde al ejemplo pequeño.}{Muestras finales de la búsqueda adaptativa con seis observaciones.}

\section{Evidencia numérica}
\Act{PARTE III}{Evidencia numérica}{Concordancia con referencias, eficiencia y sensibilidad.}
\FigureSlide{14_selected_s_agreement}{Concordancia con la referencia de malla densa}{En 1,920 superficies sintéticas y 384 financieras, los valores seleccionados de $S$ coinciden estrechamente con la referencia densa.}{Experimentos congelados; una malla finita no certifica todas las raíces.}

\FigureSlide{15_evaluation_efficiency}{Cuántas evaluaciones se ahorran}{El método adaptativo necesitó 78.35 evaluaciones en promedio frente a 5,001 en superficies sintéticas, y 91.89 frente a 5,001 en financieras.}{Resultados congelados; cada evaluación adaptativa también devuelve derivadas.}

\FigureSlide{16_selection_error_distributions}{Magnitud de las discrepancias numéricas}{La distribución de $|S^*_{\mathrm{adaptive}}-S^*_{\mathrm{dense}}|$ y del arrepentimiento complementa los porcentajes de recuperación.}{La referencia es una malla finita; el umbral de visualización no es un cero matemático.}

\FigureSlide{17_cost_by_order_horizon_window}{Cómo varía el costo con el orden y el horizonte}{El mapa de calor describe evaluaciones según el orden $d$, la ventana $L$ y el horizonte $h$. No constituye una cota de complejidad.}{Cuadrícula sintética congelada de 1,920 superficies.}

\FigureSlide{18_sensitivity_tradeoff}{Sensibilidad y costo computacional}{El refinamiento de extremos es la protección más clara. Otras variaciones cambian el costo sin justificar reajustar los parámetros previamente fijados.}{Análisis de un factor por vez sobre 216 superficies.}

\FigureSlide{19_exact_endpoint_frequencies}{Con qué frecuencia ganan los extremos}{En el estudio sintético, $S=0$ se seleccionó 312 veces y $S=1$ 34. En el financiero, las cifras fueron 53 y 10.}{Selección de objetivos numéricos, no comparación predictiva externa.}

\FigureSlide{20_adversarial_method_comparison}{Recuperación frente a casos analíticos conocidos}{La búsqueda adaptativa en $S$ recuperó 240/240 óptimos pertinentes; la búsqueda uniforme en $\log\lambda$, 226/240. Las inflexiones estacionarias son otro problema.}{Banco adversarial con soluciones conocidas por construcción.}

\FigureSlide{21_epsilon_candidate_spacing}{Por qué separar candidatos cercanos}{La distancia $\varepsilon$ suprime mínimos secundarios próximos solo después de detectarlos y ordenarlos. Se conserva el mejor mínimo interior encontrado.}{Resumen congelado de candidatos; no altera la optimización.}

\section{Interpretación aplicada}
\Act{PARTE IV}{Interpretación aplicada}{Qué significan los distintos mínimos para tendencias y pronósticos.}
\FigureSlide{22_gdp_unimodal_contrast}{PIB: la multimodalidad no es universal}{Ninguna de las 64 configuraciones predefinidas de PIB tuvo dos mínimos representativos separados por $\varepsilon$. Este caso sirve de contraste.}{Caso de PIB elegido con desarrollo; la prueba es diagnóstica.}

\FigureSlide{23_spy_rank_reversal}{SPY: la clasificación puede invertirse en prueba}{La validación favoreció $S=0.95639$ sobre $S=0$, mientras que el bloque reservado favoreció al segundo candidato. La prueba no intervino en la selección.}{Ejemplo congelado de SPY; decisiones tomadas solo con desarrollo.}

\FigureSlide{24_aapl_separated_minima}{AAPL: dos mínimos pueden ser poco competitivos}{Los candidatos $S\approx0.629$ y $S\approx0.997$ están muy separados; el segundo tiene errores mayores en validación y prueba.}{Ejemplo congelado de AAPL.}

\FigureSlide{25_btc_three_minima}{Bitcoin: tres regímenes numéricos}{En el caso elegido hay mínimos cerca de $S\approx0.845$, $0.029$ y $0.998$. La multiplicidad no implica que todos pronostiquen bien.}{Ejemplo congelado de BTC; los resultados de prueba son retrospectivos.}

\FigureSlide{26_validation_test_rank_comparison}{Separación de mínimos y precisión predictiva}{En SPY se invierte la clasificación; en BTC el segundo candidato tiene un error de prueba aproximadamente 14.7 veces mayor que el ganador de CV.}{Resultados congelados de candidatos; MSE de precios logarítmicos.}

\FigureSlide{27_multimodality_by_asset}{Frecuencia de varios mínimos en la cuadrícula estudiada}{De 64 configuraciones por serie, tienen varios candidatos: PIB 0, SPY 7, AAPL 10 y BTC 9. No representa una frecuencia general del mercado.}{Exploración de desarrollo congelada; separación $\varepsilon=0.10$.}

\FigureSlide{28_actual_temporal_splits}{Cómo evitamos utilizar el bloque de prueba}{Los orígenes de entrenamiento y validación preceden al bloque reservado. La selección de configuración y candidatos se realizó sin consultar ese bloque.}{Fechas congeladas del protocolo aplicado a SPY.}

\section{Conclusiones}
\begin{frame}{¿Qué resultados hemos obtenido?}
\begin{enumerate}
\setlength{\itemsep}{.45em}
\item La función objetivo de validación de pronósticos puede ser \textbf{multimodal}; ambos extremos representan modelos bien definidos.
\item La búsqueda adaptativa recuperó los mínimos de las referencias finitas probadas utilizando solo \textbf{1.57--1.84\%} de las evaluaciones de la malla densa.
\item Una construcción basada en \textbf{Sturm} cuenta exactamente las raíces estacionarias positivas de una instancia racional pequeña; no certifica casos grandes.
\item Diferentes mínimos pueden producir \textbf{tendencias y pronósticos distintos}; el ganador de validación no siempre gana en prueba.
\end{enumerate}
\SourceLine{Los resultados corresponden a configuraciones, datos y protocolos de comparación documentados.}
\end{frame}

\begin{frame}{Límites de las conclusiones}
\begin{columns}[T,totalwidth=\linewidth]
\begin{column}{.49\linewidth}
\begin{block}{Resultados sustentados}
\begin{itemize}
\item Derivadas analíticas para el modelo de extrapolación fijado.
\item Recuperación empírica frente a mallas densas en los casos probados.
\item Evaluación exacta de los modelos de frontera.
\item Conteo de raíces por Sturm en instancias racionales finitas.
\end{itemize}
\end{block}
\end{column}
\begin{column}{.49\linewidth}
\begin{block}{Resultados no demostrados}
\begin{itemize}
\item Recuperación universal por muestreo adaptativo finito.
\item Certificación global para ventanas de longitud arbitraria.
\item Mejora uniforme de precisión predictiva externa.
\item Predictibilidad financiera o utilidad para operar.
\end{itemize}
\end{block}
\end{column}
\end{columns}
\SourceLine{La concordancia con una malla densa no prueba la detección exhaustiva de raíces.}
\end{frame}

\begin{frame}{Reproducibilidad y referencias}
\textbf{Documento de trabajo:} \emph{Recuperación numérica de mínimos multimodales de suavizamiento óptimo para pronóstico: búsqueda adaptativa y aislamiento de Sturm}.

\medskip
\textbf{Código y experimentos:}\\
\url{https://github.com/heri-espino/Trend-Estimation}

\medskip
\textbf{Generación de figuras:}\\
\texttt{presentation/build\_figures.py} lee resultados congelados y genera ejemplos pedagógicos.\\
\texttt{presentation/slides.tex} integra los 28 archivos PDF generados.

\medskip
\textbf{Referencias principales:} Whittaker (1923); Guerrero (2007, 2008); Craven \& Wahba (1979); Hart (1994); Weinert (2007); Biessy (2026); Brent (1971). Teoría de Sturm según Basu et al. (2006).

\SourceLine{La bibliografía completa se encuentra en bibliografia.bib del artículo numérico.}
\end{frame}
\end{document}
